% Zdroj: PDF priklady-podzim-2025.pdf, fyzická str. 1. Značka: společné okruhy. Zařazeno: I2.
\section{Třídy složitosti}
\szzotazka{podzim 2025}{1}{Třídy složitosti (P/NP, redukce, NP-úplnost)}

\noindent\textbf{Zadání.}\par
\begin{enumerate}
\item Definujte \textit{třídy složitosti} \textbf{P} a \textbf{NP}.
\item Definujte \textit{převod} (redukci) mezi rozhodovacími problémy (jazyky).
\item Definujte \textbf{NP}-\textit{úplnost} rozhodovacího problému.
\item Rozhodněte, zda následující problém leží v~třídě \textbf{NP}, a~odpověď zdůvodněte:
  Je dáno přirozené číslo $x$ zapsané v~desítkové soustavě. Existují přirozená čísla $a$
  a~$b$ taková, že $1 < a, b < x$, $x = a \cdot b$ a~desítkové zápisy $a$ i~$b$ bez
  počátečních nul jsou stejně dlouhé?
\end{enumerate}

\medskip
\noindent\textbf{Náčrt řešení.}\par
Definice viz Průvodce labyrintem algoritmů, kapitoly 19.1 a~19.3. Problém leží v~třídě
\textbf{NP}, jako certifikát stačí použít číslo $a$. Verifikátor pak vypočte $b = x/a$
(pokud by při dělení vyšel zbytek, zamítne) a~zkontroluje, že $1 < a, b < x$ a~desítkové
zápisy $a$ a~$b$ jsou stejně dlouhé. Délka certifikátu je lineární a~časová složitost
verifikátoru polynomiální, obojí vzhledem k~délce vstupu (počtu číslic $x$).
